% working-notes/01-conventions.tex
% Project-specific conventions and house-style rules.
% The scaffold skill pre-fills this from the mark scheme + your context.

\begin{notesbox}[Project conventions]
\small

\textbf{Citation style.} \todo{numeric / author-year / other}.

\textbf{Units.} \todo{SI via \texttt{\textbackslash SI\{value\}\{unit\}} from siunitx, British spelling for unit names.}

\textbf{Tables and figures.} Every table and figure introduced in prose
\emph{before} it appears. Use \texttt{\textbackslash cref\{label\}} for
cross-references, never hardcode "Table~2" in prose.

\textbf{Numerical precision.} \todo{e.g.\ 3 sig figs for headline numbers.}

\textbf{Marker words.} Use \texttt{\textbackslash qContext},
\texttt{\textbackslash qOptions}, \texttt{\textbackslash qDecision},
\texttt{\textbackslash qJustification}, \texttt{\textbackslash qFindings} for
structured reasoning prose.

\textbf{Drafting markers.} \texttt{\textbackslash todo\{...\}},
\texttt{\textbackslash needcite\{description\}}, \texttt{\textbackslash gap\{description\}}
all strip-greppable before submission.

\textbf{Forbidden phrases.} No "additionally", "furthermore", "moreover",
"it is worth noting", "we explored", "interestingly", "it can be seen
that", "the analysis reveals". The marker word carries the rhetorical turn.

\textbf{Register.} Third-person or agentive. No first-person narration
of the thinking process.

\end{notesbox}
